% ============================================================================
%  CAPÍTULO II: OBJETIVOS
% ============================================================================

\section{Objetivos}

\subsection{Objetivo General}

Evaluar técnica y arquitectónicamente las principales herramientas de código abierto (*Open Source*) de monitoreo para la defensa, con el propósito de formular una propuesta de ubicación óptima en una red empresarial y validar experimentalmente las capacidades de detección pasiva (IDS) y prevención activa (IPS) sobre una infraestructura GNU/Linux.

\subsection{Objetivos Específicos}

\begin{enumerate}
  \item \textbf{Efectuar un análisis comparativo técnico y riguroso} entre dos herramientas líderes de la industria (\textbf{Suricata} y \textbf{Snort 3}), examinando sus modelos de concurrencia y paralelismo multihilo, motores de coincidencia de patrones con aceleración por hardware, mecanismos de integración con el kernel Linux y soporte de protocolos empresariales críticos.
  \item \textbf{Diseñar una propuesta arquitectónica de seguridad} que defina con claridad la ubicación estratégica de los monitores y sensores dentro de la topología de red de una organización, justificando el balance entre visibilidad perimetral (tráfico Norte-Sur) y visibilidad interna de movimientos laterales (tráfico Este-Oeste).
  \item \textbf{Implementar y verificar experimentalmente} el motor de seguridad seleccionado en la infraestructura del servidor central (\texttt{dc1}, IP \texttt{192.168.0.2}), demostrando mediante pruebas controladas de laboratorio la generación de alertas en modo IDS y el descarte efectivo de paquetes en tiempo real en modo IPS.
  \item \textbf{Documentar los procedimientos operativos y las reglas de seguridad} personalizadas, estructurando un manual técnico reproducible bajo normativas académicas y estándares internacionales de ciberseguridad.
\end{enumerate}
